\subsection{Translation Semantics and Prover Encodings}

A central challenge in multi-prover verification is ensuring that the property being proved is semantically preserved across fundamentally different logical systems. In this section, we formally specify the atomicity properties for each framework and detail their concrete encoding in SymPy, Z3, Lean~4, and Coq. We conclude with a critical discussion of the trust assumptions inherent in our translation layer.

\subsubsection{Formal Specification of Multi-Axial Atomicity}

We first restate the four atomicity criteria from the \Apeiron{} Framework in precise mathematical language, which serves as the \textit{specification} that all prover backends must implement.

\begin{enumerate}[label=(\roman*), leftmargin=*]
    \item \textbf{Boolean Atomicity:} For an observable $x$, let $B$ be a Boolean algebra with partial order $\leq$. $x$ is a Boolean atom iff $x \neq \bot$ and there exists no $y \in B$ such that $\bot < y < x$.
    \item \textbf{Measure-Theoretic Atomicity:} For a set $A$ and a measure $\mu$, $A$ is a measure atom iff $\mu(A) > 0$ and for every measurable $B \subseteq A$, either $\mu(B) = 0$ or $\mu(B) = \mu(A)$.
    \item \textbf{Categorical Atomicity:} For an object $X$ in a category $\mathcal{C}$, $X$ is a zero object (and thus atomic) iff it is both initial (unique morphism to every object) and terminal (unique morphism from every object).
    \item \textbf{Information-Theoretic Atomicity:} For a string $s$, let $K(s)$ denote its Kolmogorov complexity. $s$ is information-theoretically atomic iff $K(s) \approx |s|$; i.e., it is algorithmically incompressible.
\end{enumerate}

In practice, these abstract definitions are instantiated on the concrete value stored in an \texttt{UltimateObservable} (e.g., an integer, a set, a string). The \texttt{AtomicityTheoremGenerator} translates these value-specific claims into prover-native representations.

\subsubsection{Encoding Across Provers}

Table~\ref{tab:prover_encodings} summarizes how each atomicity framework is encoded for the four supported provers. The encodings reflect the unique strengths and idiomatic constructs of each system.

\begin{table}[htbp]
\centering
\caption{Concrete encodings of atomicity properties for each prover.}
\label{tab:prover_encodings}
\small
\begin{tabular}{p{2.2cm} p{2.8cm} p{2.8cm} p{2.5cm} p{2.5cm}}
\toprule
\textbf{Framework} & \textbf{SymPy} & \textbf{Z3} & \textbf{Lean 4} & \textbf{Coq} \\
\midrule
Boolean (int) & \texttt{And(atomic, Not(decomposable))} & \texttt{Not(Exists(x, And(x>1, x<val, val\%x==0)))} & \texttt{isAtom val} (via divisor search) & \texttt{is\_prime} (for primes) \\
Measure (set) & \texttt{len(set)==1} & \texttt{Length(set) == 1} & \texttt{cardinality} equality & \texttt{length} property \\
Category & \texttt{And(initial, terminal)} & \texttt{initial \&\& terminal} & Zero object construction & Trivial (placeholder) \\
Information & \texttt{compress\_ratio >= 0.85} & \texttt{compressed\_len >= 0.85 * original\_len} & Trivial (placeholder) & Trivial (placeholder) \\
\bottomrule
\end{tabular}
\end{table}

\paragraph{SymPy.} The generator constructs Boolean expressions using SymPy's symbolic primitives (\texttt{Symbol}, \texttt{And}, \texttt{Not}, \texttt{Implies}). For information-theoretic atomicity, it computes the zlib compression ratio and creates an expression asserting the ratio exceeds a threshold. Verification succeeds if \texttt{simplify\_logic} does not reduce the formula to \texttt{False}.

\paragraph{Z3.} Z3's SMT-LIB language supports quantifiers and arithmetic, enabling direct encoding of the existential conditions. For Boolean atomicity, the generator emits a formula asserting the non-existence of a proper divisor. For measure atomicity, it asserts that the set has exactly one element. The solver checks for satisfiability of the negated claim; \texttt{unsat} indicates a successful proof.

\paragraph{Lean 4 and Coq.} For these interactive proof assistants, the generator produces a theorem statement and a proof script. Currently, non-trivial proofs are only generated for Boolean atomicity of integers; for other frameworks, a trivial proof (\texttt{trivial} in Lean, \texttt{Proof. trivial. Qed.} in Coq) is used as a placeholder. Fully fleshed-out formalizations of measure and category theory in Lean's mathlib or Coq's standard library are planned for future work.

\subsubsection{Limitations of the Translation Layer}

The current implementation employs \textit{string templating} to generate prover-specific code. While pragmatic, this approach introduces several trust and correctness concerns:

\begin{itemize}
    \item \textbf{No formal guarantee of semantic preservation:} There is no proof that the generated Lean theorem \texttt{isAtom 1} corresponds exactly to the Python function \texttt{boolean\_atomicity(1)}. The translation relies on the programmer's manual alignment of definitions.
    \item \textbf{Syntactic fragility:} Observable IDs containing special characters (e.g., backticks, quotes) can break the generated Lean or Coq syntax. Basic sanitization is performed, but a robust parser is lacking.
    \item \textbf{Incomplete coverage:} Non-Boolean frameworks currently use placeholder proofs in Lean and Coq, meaning that cross-verification for those frameworks is only as strong as SymPy and Z3.
    \item \textbf{Lack of error recovery:} If a prover fails due to a malformed script, the system simply marks that prover as unsuccessful; it cannot diagnose or correct the translation error.
\end{itemize}

These limitations are intrinsic to the current rapid-prototyping stage of the \Apeiron{} verification module. For production-grade formal assurance, the translation layer should be rebuilt around an \textbf{Abstract Syntax Tree (AST)} in Python, with formally verified serializers for each target language. The multi-prover consensus mechanism described in Section~3 provides a pragmatic, albeit heuristic, safeguard against individual translation errors.

In the next section, we detail the theorem generation process, which builds upon the encodings described here.